\scanpage{5}
(\deleted{\unsure{normalement}}\inserted{en pratique}, $\widetilde T$ est un espace vectoriel réel de dim finie, $\pi$ un $\mathbb Z$-module libre de r. f.). \deleted{On a alors, pour $x,y\in$}\inserted{Soit $P$ un $U$-torseur,} qui hérite donc d'une topologie, d'où un groupoïde fondamental. On va décrire de façon purement algébrique ce groupoïde. On a \unsure{en effet}, pour $x,y\in P$, un isom \deleted{\unsure{de}} canonique
\[
(7)\qquad \Pi(\unsure{x|a},y)=\widetilde U_{y-x}
\qquad(\text{fibre de }\widetilde U\text{ en }y-x\in U)
\]
La composition des flèches dans $\Pi_1P$ étant donnée par l'addition dans $\widetilde U$.

\deleted{Soit} Cette formule décrit \deleted{\unsure{donc}}\inserted{un groupoïde}, en fait, pour \unsure{tte} extension de groupes abéliens
\[
\deleted{\unsure{pour}}\quad(8)\qquad0\longrightarrow\pi\longrightarrow\mathcal E\longrightarrow\gamma\longrightarrow0
\]
et un torseur $\varepsilon$ sous $\gamma$ -- on trouve un \inserted{$\pi$-}groupoïde \inserted{connexe} $\mathcal E(\varepsilon)$, ayant $\varepsilon$ comme ens de ses objets, et ses ens. de flèches donnés par
\[
\operatorname{Hom}(x,y)=\mathcal E_{y-x}
\]
\margin{\unsure{marche dès que} \unsure{l’on a un} hom de groupes $\mathcal E\longrightarrow\gamma$, \unsure{(comm. ?)} pas néc. surjectif, la surjectivité \unsure{(9)|($\varphi$)} \unsure{assure} la connexité de $\mathcal E(\varepsilon)$.}
la composition des flèches étant donnée par l'addition dans $\mathcal E$. On peut aussi décrire ce groupoïde par l'ens des flèches
\[
(10)\qquad \operatorname{Fl}\mathcal E(\varepsilon)=\deleted{\unsure{\mathcal E\times\varepsilon}}\mathcal E\times\varepsilon
\]
\mspage{4}
